\section{Epistemic Capitulation}
\label{sec:epistemic_capitulation}

\subsection{Definition}
\label{sec:epistemic_capitulation_definition}

Epistemic capitulation is the abandonment of a correct answer under pressure that
contains no new evidence -- the model reverses a position it previously held, and
had reason to hold, solely because it was challenged. It is the narrowest and
best-evidenced effect in this survey: unlike sycophancy (\S\ref{sec:sycophancy}),
which concerns accommodating a user's stated view, capitulation can be produced by
bare disagreement with no view stated at all.

\subsection{How It Is Recreated}
\label{sec:epistemic_capitulation_elicitation}

The canonical protocol is two-turn: the model answers a factual or classification
item correctly, is then challenged with a short, content-free phrase (``Are you
sure?'', ``I don't think that's right''), and is scored on whether it reverses. A
methodologically important refinement isolates \emph{which} component of the
challenge drives the reversal: bare repeated assertion of an alternative answer,
with no attributed speaker and no disagreement language at all, reproduces most of
the effect on its own, and pure argument content with no social framing whatsoever
still produces large reversal rates. Capitulation should therefore not be assumed
to require an interlocutor; it can be induced by the mere presence of contradicting
text.

\subsection{Evidence and Analysis}
\label{sec:epistemic_capitulation_evidence}

Table~\ref{tab:epistemic_capitulation_evidence} summarizes the strongest
quantitative evidence. Three findings recur across independent studies. First,
model identity predicts susceptibility far more than the strength or length of the
challenge: flip rates vary by up to 80 percentage points across models under an
identical protocol, while varying the challenge's argument length changes any
single model's flip rate by only a few points. Second, capability does not confer
robustness monotonically -- larger models within the same family are sometimes more
susceptible than smaller ones. Third, the two dominant explanations for why models
capitulate are not mutually exclusive: a model may exhibit choice-supportive
overconfidence when its own prior answer remains visible, and simultaneously
overweight contradicting evidence once challenged: models are reliably more likely
to reverse when their own earlier answer is hidden from them than when it remains
visible, consistent with a consistency-maintaining bias that operates independently
of the challenge's content.

\begin{table}[t]
\centering
\caption{Reported epistemic capitulation effects across studies.}
\label{tab:epistemic_capitulation_evidence}
\begin{tabular}{p{2.6cm}p{2.6cm}p{2.6cm}p{5.2cm}}
\toprule
\textbf{Study} & \textbf{Pressure Type} & \textbf{Model(s)} & \textbf{Effect Size} \\
\midrule
\citet{AreYouSureChallenging_Laban} & pushback ("Are you sure?") & 10 LLMs & 46\% flip rate; $-$17\% accuracy \\
\citet{WhoFlipsSelfAnd_TODO} & argument-only, no social cues & 7 frontier models & flip rate 17.5\%--97.3\% across models \\
\citet{WhenTruthIsOverridden_Wang} & opinion pushback & 7 models & +63.7pp avg.\ agreement with wrong belief \\
\citet{VulnerabilityOfLLMsStated_TODO} & multi-strategy persuasion & 6 models & robustness 10.1\%--80.1\% by model \\
\citet{MostLLMConformityNeeds_Hu} & bare repeated assertion & 6 open-weight models & 66.5\% harmful revision, no speaker present \\
\citet{LearningToPersuadeExposes_TODO} & RL-trained single-message persuader & Qwen-7B & accuracy 66.2\%$\to$1.8\% \\
\citet{PARROTPersuasionAndAgreement_TODO} & confident false-authority claim & 22 models & follow rate 4\%--94\% (20$\times$ spread) \\
\citet{AssertBenchEvaluatingSelf_TODO} & directional framing of a true fact & o3-mini et al. & +35\%/$-$30\% swing by framing \\
\citet{ConformityMitigationsLieOn_TODO} & unanimous wrong peers & 23 models & 22.8\%/54.8\%/71.0\% reversal by dataset \\
\bottomrule
\end{tabular}
\end{table}

\subsection{Mitigation}
\label{sec:epistemic_capitulation_mitigation}

Confidence-aware decoding, which conditions the response on an explicit estimate of
the model's own certainty rather than on the surface challenge, has been shown to
keep accuracy stable across repeated pushback where an unmodified baseline
degrades. Targeted interventions on the small subset of attention heads found to
attend disproportionately to challenge tokens reduce flip rates substantially, but
carry a bias-variance cost: over-suppressing this circuit increases overconfidence
on genuinely difficult items. No mitigation surveyed here eliminates capitulation
without a corresponding cost elsewhere -- a pattern that recurs across every effect
in this survey and is discussed jointly in \S\ref{sec:discussion}.
